\documentclass[11pt,a4paper]{article}

% Packages
\usepackage[margin=1in]{geometry}
\usepackage{amsmath}
\usepackage{amssymb}
\usepackage{graphicx}
\usepackage{listings}
\usepackage{xcolor}
\usepackage{hyperref}
\usepackage{booktabs}
\usepackage{fancyhdr}
\usepackage{tcolorbox}
\usepackage{enumitem}

% R code formatting
\lstset{
    language=R,
    basicstyle=\ttfamily\small,
    keywordstyle=\color{blue},
    commentstyle=\color{gray},
    stringstyle=\color{red},
    numbers=left,
    numberstyle=\tiny\color{gray},
    stepnumber=1,
    numbersep=5pt,
    backgroundcolor=\color{white},
    showspaces=false,
    showstringspaces=false,
    showtabs=false,
    frame=single,
    tabsize=2,
    captionpos=b,
    breaklines=true,
    breakatwhitespace=false,
    escapeinside={(*@}{@*)},
    xleftmargin=2em,
    framexleftmargin=1.5em
}

% Custom boxes
\tcbuselibrary{skins,breakable}
\newtcolorbox{keypoint}{
    colback=blue!5!white,
    colframe=blue!75!black,
    fonttitle=\bfseries,
    title=Key Point,
    breakable
}

\newtcolorbox{warning}{
    colback=red!5!white,
    colframe=red!75!black,
    fonttitle=\bfseries,
    title=Common Mistake,
    breakable
}

\newtcolorbox{rcode}{
    colback=green!5!white,
    colframe=green!75!black,
    fonttitle=\bfseries,
    title=R Code Example,
    breakable
}

\newtcolorbox{example}{
    colback=yellow!5!white,
    colframe=orange!75!black,
    fonttitle=\bfseries,
    title=Example,
    breakable
}

% Header and footer
\pagestyle{fancy}
\fancyhf{}
\rhead{Mixed Effects Models Guide}
\lhead{TAD Analysis}
\rfoot{Page \thepage}

% Title
\title{\textbf{ Guide to Mixed Effects Models} \\ \large Understanding and Interpreting TAD Persistence and Reemergence Analysis}
\author{Prepared for Chapter 2 Analysis}
\date{\today}

\begin{document}

\maketitle
\tableofcontents
\newpage

%==============================================================================
\section{Introduction: What You Need to Know}
%==============================================================================

This guide will help you understand, run, and interpret mixed effects models for your TAD (Transboundary Animal Disease) analysis. By the end, you'll be able to:

\begin{itemize}[leftmargin=*]
    \item Understand what mixed effects models do and why you need them
    \item Interpret odds ratios (OR) and hazard ratios (HR)
    \item Map statistical output to your five hypotheses
    \item Run the analysis in R with confidence
    \item Write clear results for your chapter
\end{itemize}

\begin{keypoint}

\begin{enumerate}
    \item Can you interpret the coefficients from mixed effects models?
    \item Do you know which coefficients must be significant to confirm each hypothesis?
\end{enumerate}
This guide answers both questions clearly.
\end{keypoint}

%==============================================================================
\section{Part 1: The Foundation - Understanding Regression}
%==============================================================================

\subsection{The Core Equation Structure}

Every regression model has this fundamental form:

\begin{equation}
\boxed{\text{OUTCOME} = \text{INTERCEPT} + \text{SLOPE} \times \text{PREDICTOR} + \text{ERROR}}
\end{equation}

Written mathematically:

\begin{equation}
Y = \beta_0 + \beta_1 \times X + \varepsilon
\end{equation}

\textbf{Where:}
\begin{itemize}[leftmargin=*]
    \item $Y$ = outcome variable (what you're trying to predict)
    \item $\beta_0$ = intercept (baseline value when $X = 0$)
    \item $\beta_1$ = slope/coefficient (effect of predictor)
    \item $X$ = predictor variable (what you think influences $Y$)
    \item $\varepsilon$ = error term (unexplained variation)
\end{itemize}

\begin{example}
\textbf{Simple Example: Height and Age}

If we're predicting height from age:
\[
\text{Height} = 80\text{ cm} + 2\text{ cm/year} \times \text{Age} + \text{error}
\]

\begin{itemize}
    \item $\beta_0 = 80$ cm (baseline height at birth)
    \item $\beta_1 = 2$ cm/year (growth rate)
    \item For a 5-year-old: Height $= 80 + 2(5) = 90$ cm
\end{itemize}
\end{example}

\subsection{The Greek Letters (Don't Fear Them!)}

Statistics uses Greek letters, but they're just symbols with specific meanings:

\begin{table}[h]
\centering
\begin{tabular}{@{}clp{8cm}@{}}
\toprule
\textbf{Symbol} & \textbf{Name} & \textbf{Meaning} \\
\midrule
$\beta$ & beta & Coefficient/slope (the effect size) \\
$\beta_0$ & beta-zero & Intercept (baseline value) \\
$\beta_1$ & beta-one & Effect of predictor 1 \\
$\varepsilon$ & epsilon & Error term (unexplained variation) \\
$\sigma^2$ & sigma-squared & Variance (how spread out data is) \\
$u$ & u & Random effect (group-specific deviation) \\
\bottomrule
\end{tabular}
\end{table}

\begin{keypoint}
In your analysis, you'll report $\beta$ values, but you'll convert them to OR or HR for interpretation. The sign of $\beta$ (positive or negative) tells you the \textbf{direction} of the effect.
\end{keypoint}

%==============================================================================
\section{Part 2: Why Mixed Effects Models?}
%==============================================================================

\subsection{The Clustering Problem}

Your TAD data has a hierarchical structure:

\begin{verbatim}
Disease: ASF
  ├── Country: Romania
  │     ├── County 1: Outbreak A, Outbreak B, Outbreak C
  │     ├── County 2: Outbreak D, Outbreak E
  │     └── County 3: Outbreak F
  └── Country: Poland
        └── County 4: Outbreak G, Outbreak H
        
Disease: FMD
  └── Country: Turkey
        └── County 5: Outbreak I, Outbreak J
\end{verbatim}

\textbf{The problem:} Outbreaks A and B (same county) are more similar to each other than Outbreak A and G (different countries). This is called \textbf{clustering}.

\subsection{What Happens If You Ignore Clustering?}

\begin{warning}
If you use simple regression and ignore clustering:
\begin{itemize}
    \item Your p-values will be \textbf{too small} (false positives)
    \item Your confidence intervals will be \textbf{too narrow} (overconfident)
    \item You'll conclude things are significant when they're not
  
\end{itemize}
\end{warning}

\subsection{How Mixed Effects Models Fix This}

Mixed effects models have two types of effects:

\begin{enumerate}
    \item \textbf{Fixed effects} ($\beta$): The average effect across all groups
        \begin{itemize}
            \item This is what you test for your hypotheses
            \item Example: ``Livestock density increases persistence''
        \end{itemize}
    
    \item \textbf{Random effects} ($u$): Group-specific deviations
        \begin{itemize}
            \item Accounts for clustering
            \item Allows each disease/country/county to have its own baseline
            \item You don't directly interpret these, but they're crucial for correct standard errors
        \end{itemize}
\end{enumerate}

\begin{keypoint}
Think of it this way:
\begin{itemize}
    \item \textbf{Fixed effects:} ``On average, what's the effect of livestock density?''
    \item \textbf{Random effects:} ``But ASF naturally persists more than FMD, Romania has different baseline risk than Poland, etc.''
\end{itemize}
\end{keypoint}

\subsection{Your Three-Level Random Effects Structure}

Your model accounts for clustering at three levels:

\begin{equation}
\text{persistence} = \beta_0 + \beta_1 X_1 + \cdots + \underbrace{u_{\text{disease}} + u_{\text{country}|\text{disease}} + u_{\text{admin2}|\text{country}|\text{disease}}}_{\text{Random effects}} + \varepsilon
\end{equation}

\textbf{In R syntax:}
\begin{lstlisting}[language=R]
(1|disease/country/admin2)
\end{lstlisting}

This means:
\begin{itemize}
    \item Each disease has its own baseline (some diseases naturally persist more)
    \item Each country within a disease has its own baseline
    \item Each admin unit within a country has its own baseline
\end{itemize}

%==============================================================================
\section{Part 3: Your Two Models}
%==============================================================================

\subsection{Model 1: Persistence (Logistic Regression)}

\textbf{Research question:} What predicts whether an outbreak recurs within 6-24 months?

\textbf{Outcome variable:} Binary (0 = no recurrence, 1 = recurrence)

\textbf{Statistical model:} Mixed-effects logistic regression

\textbf{The equation:}
\begin{equation}
\log\left(\frac{P(\text{persistence})}{1-P(\text{persistence})}\right) = \beta_0 + \beta_1 X_1 + \beta_2 X_2 + \cdots + u_{\text{groups}} + \varepsilon
\end{equation}

\textbf{Output:} Odds Ratio (OR)

\begin{rcode}
\begin{lstlisting}[language=R]
# Load required package
library(lme4)

# Fit persistence model
model_persist <- glmer(
  persistence ~ 
    livestock_density + 
    vaccination_coverage + 
    wildlife_proximity + 
    tree_cover + 
    water_bodies + 
    road_density + 
    (1|disease/country/admin2),  # Random effects
  family = binomial,  # For binary outcome
  data = dat
)

# View results
summary(model_persist)

# Get odds ratios (exponentiate coefficients)
exp(fixef(model_persist))

# Get confidence intervals
exp(confint(model_persist))
\end{lstlisting}
\end{rcode}

\subsection{Model 2: Reemergence (Cox Proportional Hazards)}

\textbf{Research question:} How long until disease returns after $\geq 2$ years absence?

\textbf{Outcome variable:} Time-to-event (months until reemergence)

\textbf{Statistical model:} Cox proportional hazards with frailty (random effects)

\textbf{The equation:}
\begin{equation}
h(t|X) = h_0(t) \times \exp(\beta_1 X_1 + \beta_2 X_2 + \cdots + u_{\text{groups}})
\end{equation}

\textbf{Output:} Hazard Ratio (HR)

\begin{rcode}
\begin{lstlisting}[language=R]
# Load required package
library(survival)

# Fit reemergence model
model_reemerge <- coxph(
  Surv(time, event) ~  # time = months, event = 1 if reemerged
    livestock_density + 
    vaccination_coverage + 
    wildlife_proximity + 
    tree_cover + 
    water_bodies + 
    road_density + 
    frailty(disease/country/admin2),  # Random effects
  data = dat
)

# View results
summary(model_reemerge)

# Hazard ratios are already in exp(coef) column
# No need to exponentiate
\end{lstlisting}
\end{rcode}

%==============================================================================
\section{Part 4: Understanding Odds Ratios (OR)}
%==============================================================================

\subsection{From Probability to Odds to Odds Ratio}

\textbf{Step 1: Probability} (you already understand this)
\begin{itemize}
    \item Probability = chance of something happening
    \item Ranges from 0 to 1
    \item Example: 30\% chance of rain = 0.30 probability
\end{itemize}

\textbf{Step 2: Odds} (slightly different)
\begin{equation}
\text{Odds} = \frac{P(\text{event})}{P(\text{no event})} = \frac{P}{1-P}
\end{equation}

\begin{example}
\textbf{Converting Probability to Odds:}
\begin{itemize}
    \item $P = 0.50$ (50\% chance) $\rightarrow$ Odds = $\frac{0.50}{0.50} = 1.0$ (``even odds'' or ``1:1'')
    \item $P = 0.75$ (75\% chance) $\rightarrow$ Odds = $\frac{0.75}{0.25} = 3.0$ (``3:1 odds'')
    \item $P = 0.25$ (25\% chance) $\rightarrow$ Odds = $\frac{0.25}{0.75} = 0.33$ (``1:3 odds'')
\end{itemize}
\end{example}

\textbf{Step 3: Odds Ratio} (comparing two groups)
\begin{equation}
\text{OR} = \frac{\text{Odds}(\text{exposed})}{\text{Odds}(\text{unexposed})}
\end{equation}

\begin{example}
\textbf{Comparing High vs. Low Livestock Density:}

\textbf{High density:} 60 outbreaks persisted, 40 did not
\[
\text{Odds}_{\text{high}} = \frac{60}{40} = 1.5
\]

\textbf{Low density:} 30 outbreaks persisted, 70 did not
\[
\text{Odds}_{\text{low}} = \frac{30}{70} = 0.43
\]

\textbf{Odds ratio:}
\[
\text{OR} = \frac{1.5}{0.43} = 3.49
\]

\textbf{Interpretation:} High-density areas have 3.49 times the odds of persistence compared to low-density areas.
\end{example}

\subsection{The Critical OR Rules (MEMORIZE THIS)}

\begin{table}[h]
\centering
\begin{tabular}{@{}clp{7cm}@{}}
\toprule
\textbf{OR Value} & \textbf{Direction} & \textbf{Interpretation} \\
\midrule
OR = 1.0 & No effect & Predictor doesn't matter \\
OR $>$ 1.0 & Risk factor & \textbf{Increases} odds of persistence \\
OR $<$ 1.0 & Protective & \textbf{Decreases} odds of persistence \\
\midrule
OR = 2.0 & Risk & Odds are doubled (100\% increase) \\
OR = 1.5 & Risk & 50\% increase in odds \\
OR = 1.2 & Risk & 20\% increase in odds \\
OR = 0.5 & Protective & 50\% reduction in odds \\
OR = 0.8 & Protective & 20\% reduction in odds \\
\bottomrule
\end{tabular}
\end{table}

\subsection{Converting OR to Percentage Change}

\begin{equation}
\boxed{\text{Percentage change} = (\text{OR} - 1) \times 100\%}
\end{equation}

\begin{example}
\textbf{Examples:}
\begin{itemize}
    \item OR = 1.68 $\rightarrow$ $(1.68 - 1) \times 100\% = 68\%$ increase
    \item OR = 0.75 $\rightarrow$ $(0.75 - 1) \times 100\% = -25\%$ (25\% decrease)
    \item OR = 2.50 $\rightarrow$ $(2.50 - 1) \times 100\% = 150\%$ increase
\end{itemize}
\end{example}

\subsection{The Relationship: $\beta$ to OR}

In your R output, you'll see coefficients ($\beta$). To get OR:

\begin{equation}
\boxed{\text{OR} = e^{\beta} = \exp(\beta)}
\end{equation}

\begin{example}
\textbf{R Output:}
\begin{verbatim}
                      Estimate  Std.Error  z value  Pr(>|z|)
livestock_density      0.523     0.182      2.87     0.004
\end{verbatim}

\textbf{Calculate OR:}
\[
\text{OR} = e^{0.523} = 1.687 \approx 1.69
\]

\textbf{Interpretation:} A 1-SD increase in livestock density is associated with 69\% higher odds of persistence (OR = 1.69, p = 0.004).
\end{example}

\begin{rcode}
\begin{lstlisting}[language=R]
# Get coefficients
beta <- fixef(model_persist)["livestock_density"]
beta  # Shows: 0.523

# Convert to OR
OR <- exp(beta)
OR  # Shows: 1.687

# Calculate percentage change
pct_change <- (OR - 1) * 100
pct_change  # Shows: 68.7%

# Get 95% confidence interval
beta_se <- sqrt(diag(vcov(model_persist)))["livestock_density"]
CI_lower <- exp(beta - 1.96 * beta_se)
CI_upper <- exp(beta + 1.96 * beta_se)

cat("OR =", round(OR, 2), 
    "95% CI: (", round(CI_lower, 2), "-", round(CI_upper, 2), ")")
# Output: OR = 1.69 95% CI: ( 1.18 - 2.41 )
\end{lstlisting}
\end{rcode}

%==============================================================================
\section{Part 5: Understanding Hazard Ratios (HR)}
%==============================================================================

\subsection{What is a Hazard?}

\textbf{Hazard} = instantaneous risk of an event occurring at time $t$, given survival up to time $t$

Think of it as: ``At this moment, what's the risk of disease returning?''

\subsection{The Hazard Ratio}

\begin{equation}
\text{HR} = \frac{\text{Hazard}(\text{exposed})}{\text{Hazard}(\text{unexposed})} = e^{\beta}
\end{equation}

\subsection{The Critical HR Rules (MEMORIZE THIS)}

\begin{table}[h]
\centering
\begin{tabular}{@{}clp{7cm}@{}}
\toprule
\textbf{HR Value} & \textbf{Direction} & \textbf{Interpretation} \\
\midrule
HR = 1.0 & No effect & Predictor doesn't matter \\
HR $>$ 1.0 & Risk factor & \textbf{Faster} reemergence (BAD) \\
HR $<$ 1.0 & Protective & \textbf{Slower} reemergence (GOOD) \\
\midrule
HR = 2.0 & Risk & Reemergence happens twice as fast \\
HR = 1.5 & Risk & 50\% faster reemergence \\
HR = 1.2 & Risk & 20\% faster reemergence \\
HR = 0.5 & Protective & Takes twice as long to reemerge \\
HR = 0.7 & Protective & 30\% slower reemergence \\
\bottomrule
\end{tabular}
\end{table}

\begin{keypoint}
\textbf{Key Difference from OR:}
\begin{itemize}
    \item OR compares \textbf{odds} at a single point in time
    \item HR compares \textbf{rates} over time
\end{itemize}

Both use the same interpretation framework: $>1$ = risk factor, $<1$ = protective.
\end{keypoint}

\begin{example}
\textbf{Vaccination and Reemergence:}

R output shows:
\begin{verbatim}
                         coef  exp(coef)  se(coef)      z  Pr(>|z|)
vaccination_coverage   -0.523      0.593     0.188  -2.78     0.005
\end{verbatim}

\textbf{Interpretation:}
\begin{itemize}
    \item $\beta = -0.523$ (negative = protective)
    \item HR = $e^{-0.523} = 0.593$
    \item Counties with high vaccination have 59.3\% the hazard of reemergence
    \item Or: 41\% slower reemergence $(1 - 0.593) \times 100\% = 40.7\%$
\end{itemize}
\end{example}

\begin{rcode}
\begin{lstlisting}[language=R]
# Cox model already gives you HR in exp(coef) column
summary(model_reemerge)

# Extract specific HR
coef_summary <- summary(model_reemerge)$coefficients
HR_vaccination <- coef_summary["vaccination_coverage", "exp(coef)"]
HR_vaccination  # Shows: 0.593

# Calculate percentage change
pct_slower <- (1 - HR_vaccination) * 100
cat("Vaccination delays reemergence by", round(pct_slower, 1), "%")
# Output: Vaccination delays reemergence by 40.7 %

# Get confidence interval (already in summary)
CI_lower <- coef_summary["vaccination_coverage", "lower .95"]
CI_upper <- coef_summary["vaccination_coverage", "upper .95"]
\end{lstlisting}
\end{rcode}

%==============================================================================
\section{Part 6: Interpreting Your Output}
%==============================================================================

\subsection{The Four-Step Interpretation Process}

For \textbf{EVERY} predictor in your model, follow these four steps:

\begin{enumerate}[leftmargin=*]
    \item \textbf{Check the SIGN of $\beta$}
    \begin{itemize}
        \item Positive $\rightarrow$ increases odds/hazard (risk factor)
        \item Negative $\rightarrow$ decreases odds/hazard (protective)
    \end{itemize}
    
    \item \textbf{Check if SIGNIFICANT}
    \begin{itemize}
        \item Is $p < 0.05$?
        \item If NO, stop here. No biological interpretation.
    \end{itemize}
    
    \item \textbf{Calculate MAGNITUDE}
    \begin{itemize}
        \item OR/HR = $e^{\beta}$
        \item Percentage change = (OR/HR - 1) $\times$ 100\%
    \end{itemize}
    
    \item \textbf{Check PRECISION}
    \begin{itemize}
        \item Does 95\% CI cross 1.0?
        \item Narrow CI = precise, Wide CI = uncertain
    \end{itemize}
\end{enumerate}

\subsection{Example: Complete Interpretation}

\textbf{Persistence Model Output:}
\begin{verbatim}
Fixed effects:
                      Estimate  Std.Error  z value  Pr(>|z|)
livestock_density      0.523     0.182      2.87     0.004
vaccination_coverage  -0.678     0.224     -3.03     0.002
wildlife_proximity     0.312     0.148      2.11     0.035
\end{verbatim}

\textbf{Step-by-step for livestock\_density:}

\begin{enumerate}
    \item \textbf{Sign:} $\beta = 0.523$ (positive) $\rightarrow$ increases persistence odds
    
    \item \textbf{Significance:} $p = 0.004 < 0.05$ $\rightarrow$ \checkmark Significant
    
    \item \textbf{Magnitude:}
    \begin{align*}
        \text{OR} &= e^{0.523} = 1.69 \\
        \text{Percentage change} &= (1.69 - 1) \times 100\% = 69\%
    \end{align*}
    
    \item \textbf{Precision:}
    \begin{align*}
        \text{CI}_{\text{lower}} &= e^{0.523 - 1.96(0.182)} = 1.18 \\
        \text{CI}_{\text{upper}} &= e^{0.523 + 1.96(0.182)} = 2.41
    \end{align*}
    Doesn't cross 1.0 $\rightarrow$ confirms significance
\end{enumerate}

\textbf{Written interpretation:}
\begin{quote}
``Counties with high livestock density (1 SD above mean) had 69\% higher odds of outbreak persistence compared to counties with average density (OR = 1.69, 95\% CI: 1.18--2.41, $p = 0.004$), supporting Hypothesis 1.''
\end{quote}

\begin{rcode}
\begin{lstlisting}[language=R]
# Complete interpretation workflow
library(lme4)

# 1. Fit model
model <- glmer(persistence ~ livestock_density + vaccination_coverage + 
               wildlife_proximity + (1|disease/country/admin2),
               family = binomial, data = dat)

# 2. Extract results
coef_table <- summary(model)$coefficients

# 3. For each predictor, check four steps
predictor <- "livestock_density"

# Step 1: Sign
beta <- coef_table[predictor, "Estimate"]
direction <- ifelse(beta > 0, "increases", "decreases")
cat("Step 1 - Sign:", direction, "persistence\n")

# Step 2: Significance
p_value <- coef_table[predictor, "Pr(>|z|)"]
is_sig <- p_value < 0.05
cat("Step 2 - Significant:", is_sig, "(p =", round(p_value, 3), ")\n")

if(is_sig) {
  # Step 3: Magnitude
  OR <- exp(beta)
  pct_change <- (OR - 1) * 100
  cat("Step 3 - Magnitude: OR =", round(OR, 2), 
      "(", round(abs(pct_change), 1), "% change)\n")
  
  # Step 4: Precision
  se <- coef_table[predictor, "Std. Error"]
  CI_lower <- exp(beta - 1.96 * se)
  CI_upper <- exp(beta + 1.96 * se)
  cat("Step 4 - Precision: 95% CI (", 
      round(CI_lower, 2), "-", round(CI_upper, 2), ")\n")
  
  # Write interpretation
  cat("\nInterpretation:\n")
  cat("Counties with high", predictor, 
      "had", round(abs(pct_change), 0), "% ",
      ifelse(beta > 0, "higher", "lower"),
      "odds of persistence (OR =", round(OR, 2), 
      ", 95% CI:", round(CI_lower, 2), "-", round(CI_upper, 2),
      ", p =", format.pval(p_value, digits = 3), ")\n")
}

# Output:
# Step 1 - Sign: increases persistence
# Step 2 - Significant: TRUE (p = 0.004)
# Step 3 - Magnitude: OR = 1.69 ( 68.7 % change)
# Step 4 - Precision: 95% CI ( 1.18 - 2.41 )
# 
# Interpretation:
# Counties with high livestock_density had 69 % higher odds of 
# persistence (OR = 1.69 , 95% CI: 1.18 - 2.41 , p = 0.004)
\end{lstlisting}
\end{rcode}

%==============================================================================
\section{Part 7: Mapping Coefficients to Your Hypotheses}
%==============================================================================

\subsection{Your Five Hypotheses}

This section shows \textbf{exactly} which coefficients must be significant and in which direction for each hypothesis to be supported.

\subsubsection{H1: Ecological and Environmental Risk Factors}

\textbf{Prediction:} Regions with high livestock density, wildlife proximity, tree cover, and water/wetlands will have higher persistence and faster reemergence.

\textbf{Decision Rule:}

\begin{table}[h]
\centering
\small
\begin{tabular}{@{}lccl@{}}
\toprule
\textbf{Predictor} & \textbf{Expected Sign} & \textbf{Threshold} & \textbf{Model} \\
\midrule
livestock\_density & $\beta > 0$ & $p < 0.05$ & Both \\
wildlife\_proximity & $\beta > 0$ & $p < 0.05$ & Both \\
tree\_cover & $\beta > 0$ & $p < 0.05$ & Both \\
water\_bodies OR wetlands & $\beta > 0$ & $p < 0.05$ & Both \\
\bottomrule
\end{tabular}
\end{table}

\textbf{H1 is SUPPORTED if:}
\begin{itemize}
    \item \textbf{ALL} four predictors are positive AND significant in persistence model
    \item \textbf{ALL} four predictors are positive AND significant in reemergence model
\end{itemize}

\begin{rcode}
\begin{lstlisting}[language=R]
# Test H1 for persistence model
H1_predictors <- c("livestock_density", "wildlife_proximity", 
                   "tree_cover", "water_bodies")

# Extract coefficients and p-values
coef_table <- summary(model_persist)$coefficients

# Check each predictor
H1_checks <- sapply(H1_predictors, function(pred) {
  beta <- coef_table[pred, "Estimate"]
  p_val <- coef_table[pred, "Pr(>|z|)"]
  
  is_positive <- beta > 0
  is_significant <- p_val < 0.05
  
  return(is_positive & is_significant)
})

# H1 supported if ALL are TRUE
H1_supported <- all(H1_checks)

cat("H1 Supported:", H1_supported, "\n")
cat("Individual checks:\n")
print(H1_checks)

# Example output:
# H1 Supported: TRUE 
# Individual checks:
#   livestock_density wildlife_proximity         tree_cover       water_bodies 
#                TRUE               TRUE               TRUE               TRUE
\end{lstlisting}
\end{rcode}

\subsubsection{H2: Infrastructure and Accessibility}

\textbf{Prediction:} Dense road networks increase persistence and accelerate reemergence.

\textbf{Decision Rule:}

\begin{itemize}
    \item Persistence: $\beta_{\text{road\_density}} > 0$ AND $p < 0.05$
    \item Reemergence: $\beta_{\text{road\_density}} > 0$ AND $p < 0.05$
\end{itemize}

\textbf{H2 is SUPPORTED if:} Both conditions are met.

\begin{rcode}
\begin{lstlisting}[language=R]
# Test H2
# For persistence
beta_roads_persist <- coef(model_persist)["road_density"]
p_roads_persist <- summary(model_persist)$coefficients["road_density", "Pr(>|z|)"]
H2_persist <- (beta_roads_persist > 0) & (p_roads_persist < 0.05)

# For reemergence
beta_roads_reemerge <- coef(model_reemerge)["road_density"]
p_roads_reemerge <- summary(model_reemerge)$coefficients["road_density", "Pr(>|z|)"]
H2_reemerge <- (beta_roads_reemerge > 0) & (p_roads_reemerge < 0.05)

# H2 supported if both are TRUE
H2_supported <- H2_persist & H2_reemerge

cat("H2 Supported:", H2_supported, "\n")
cat("  Persistence:", H2_persist, "\n")
cat("  Reemergence:", H2_reemerge, "\n")
\end{lstlisting}
\end{rcode}

\subsubsection{H3: Control Interventions and Institutional Capacity}

\textbf{Prediction:} Vaccination and culling reduce persistence and delay reemergence, with stronger effects for persistence.

\textbf{Decision Rule - Part 1 (Basic effect):}

\begin{table}[h]
\centering
\small
\begin{tabular}{@{}lcc@{}}
\toprule
\textbf{Intervention} & \textbf{Persistence} & \textbf{Reemergence} \\
\midrule
vaccination\_coverage & $\beta < 0$, $p < 0.05$ & $\beta < 0$, $p < 0.05$ \\
culling\_coverage & $\beta < 0$, $p < 0.05$ & $\beta < 0$, $p < 0.05$ \\
\bottomrule
\end{tabular}
\end{table}

\textbf{Decision Rule - Part 2 (Stronger for persistence):}
\begin{itemize}
    \item $|\beta_{\text{persistence}}| > |\beta_{\text{reemergence}}|$ for both interventions
\end{itemize}

\begin{rcode}
\begin{lstlisting}[language=R]
# Test H3 - Part 1: Basic protective effect
interventions <- c("vaccination_coverage", "culling_coverage")

# Persistence
H3_persist <- sapply(interventions, function(var) {
  beta <- coef(model_persist)[var]
  p_val <- summary(model_persist)$coefficients[var, "Pr(>|z|)"]
  return((beta < 0) & (p_val < 0.05))
})

# Reemergence
H3_reemerge <- sapply(interventions, function(var) {
  beta <- coef(model_reemerge)[var]
  p_val <- summary(model_reemerge)$coefficients[var, "Pr(>|z|)"]
  return((beta < 0) & (p_val < 0.05))
})

H3_part1 <- all(H3_persist) & all(H3_reemerge)

# Part 2: Compare magnitudes
beta_vacc_persist <- abs(coef(model_persist)["vaccination_coverage"])
beta_vacc_reemerge <- abs(coef(model_reemerge)["vaccination_coverage"])
beta_cull_persist <- abs(coef(model_persist)["culling_coverage"])
beta_cull_reemerge <- abs(coef(model_reemerge)["culling_coverage"])

H3_part2 <- (beta_vacc_persist > beta_vacc_reemerge) & 
            (beta_cull_persist > beta_cull_reemerge)

# Full H3
H3_supported <- H3_part1 & H3_part2

cat("H3 Supported:", H3_supported, "\n")
cat("  Part 1 (protective effects):", H3_part1, "\n")
cat("  Part 2 (stronger for persistence):", H3_part2, "\n")

# Calculate effect sizes for reporting
OR_vacc_persist <- exp(-beta_vacc_persist)
HR_vacc_reemerge <- exp(-beta_vacc_reemerge)
reduction_persist <- (1 - OR_vacc_persist) * 100
reduction_reemerge <- (1 - HR_vacc_reemerge) * 100

cat("\nVaccination effect sizes:\n")
cat("  Persistence:", round(reduction_persist, 1), "% reduction\n")
cat("  Reemergence:", round(reduction_reemerge, 1), "% reduction\n")
\end{lstlisting}
\end{rcode}

\subsubsection{H4: Wildlife-Livestock Interface}

\textbf{Prediction:} Wildlife-origin outbreaks, proximity to wildlife, and forest cover interactions increase persistence and accelerate reemergence.

\textbf{Decision Rule:}

\begin{table}[h]
\centering
\small
\begin{tabular}{@{}lcc@{}}
\toprule
\textbf{Predictor} & \textbf{Expected Sign} & \textbf{Threshold} \\
\midrule
wildlife\_origin & $\beta > 0$ & $p < 0.05$ \\
distance\_to\_wildlife & $\beta < 0$ & $p < 0.05$ \\
wildlife $\times$ tree\_cover & $\beta > 0$ & $p < 0.05$ \\
\bottomrule
\end{tabular}
\end{table}

\textbf{H4 is SUPPORTED if:} All three components are significant with predicted signs in both models.

\begin{rcode}
\begin{lstlisting}[language=R]
# Test H4
H4_components <- c("wildlife_origin", 
                   "distance_to_wildlife", 
                   "wildlife_proximity:tree_cover")

# Expected signs
expected_signs <- c(1, -1, 1)  # positive, negative, positive

# Function to test H4 for one model
test_H4 <- function(model, model_type) {
  coef_table <- summary(model)$coefficients
  
  checks <- sapply(1:length(H4_components), function(i) {
    var <- H4_components[i]
    expected <- expected_signs[i]
    
    beta <- coef_table[var, "Estimate"]
    p_val <- coef_table[var, "Pr(>|z|)"]
    
    # Check sign matches expectation
    sign_correct <- (expected > 0 & beta > 0) | (expected < 0 & beta < 0)
    is_sig <- p_val < 0.05
    
    return(sign_correct & is_sig)
  })
  
  names(checks) <- H4_components
  return(checks)
}

# Test both models
H4_persist_checks <- test_H4(model_persist, "persistence")
H4_reemerge_checks <- test_H4(model_reemerge, "reemergence")

H4_supported <- all(H4_persist_checks) & all(H4_reemerge_checks)

cat("H4 Supported:", H4_supported, "\n")
cat("\nPersistence model:\n")
print(H4_persist_checks)
cat("\nReemergence model:\n")
print(H4_reemerge_checks)
\end{lstlisting}
\end{rcode}

\subsubsection{H5: Disease-Specific Risk Factor Profiles}

\textbf{Prediction:} Risk factors differ by transmission mode (vector-borne vs. direct transmission).

\textbf{Specific predictions:}
\begin{itemize}
    \item Vector-borne diseases: stronger wetland and climate effects
    \item Direct transmission diseases: stronger livestock density and intervention effects
\end{itemize}

\textbf{Testing approach:} Run disease-specific models, extract coefficients, test for heterogeneity using meta-regression.

\begin{rcode}
\begin{lstlisting}[language=R]
# H5: Disease-specific models
library(metafor)

# 1. Run disease-specific models
diseases <- unique(dat$disease)
disease_results <- data.frame()

for(disease in diseases) {
  # Subset data
  dat_sub <- subset(dat, disease == disease)
  
  # Fit model (without disease random effect)
  model_disease <- glmer(
    persistence ~ livestock_density + vaccination_coverage + 
                  wetlands + (1|country/admin2),
    family = binomial,
    data = dat_sub
  )
  
  # Extract coefficients
  coef_table <- summary(model_disease)$coefficients
  
  disease_results <- rbind(disease_results, data.frame(
    disease = disease,
    transmission = ifelse(disease %in% c("RVF", "BT", "LSD"), 
                         "vector", "direct"),
    beta_wetlands = coef_table["wetlands", "Estimate"],
    se_wetlands = coef_table["wetlands", "Std. Error"],
    beta_livestock = coef_table["livestock_density", "Estimate"],
    se_livestock = coef_table["livestock_density", "Std. Error"],
    beta_vaccination = coef_table["vaccination_coverage", "Estimate"],
    se_vaccination = coef_table["vaccination_coverage", "Std. Error"]
  ))
}

# 2. Meta-regression: Does transmission mode moderate effects?
# Test wetlands effect
meta_wetlands <- rma(yi = beta_wetlands, sei = se_wetlands, 
                     mods = ~ transmission, 
                     data = disease_results)

# Test livestock effect
meta_livestock <- rma(yi = beta_livestock, sei = se_livestock, 
                      mods = ~ transmission, 
                      data = disease_results)

# 3. Check heterogeneity
cat("Wetlands effect by transmission mode:\n")
print(meta_wetlands)
cat("\nHeterogeneity p-value:", meta_wetlands$QMp, "\n")

cat("\nLivestock effect by transmission mode:\n")
print(meta_livestock)
cat("\nHeterogeneity p-value:", meta_livestock$QMp, "\n")

# H5 supported if:
# - Wetlands heterogeneity p < 0.05 AND stronger for vector-borne
# - Livestock heterogeneity p < 0.05 AND stronger for direct
H5_wetlands <- (meta_wetlands$QMp < 0.05) & 
               (meta_wetlands$beta[2] > 0)  # moderator coefficient
H5_livestock <- (meta_livestock$QMp < 0.05) & 
                (meta_livestock$beta[2] < 0)  # direct has lower effect

H5_supported <- H5_wetlands | H5_livestock  # Either or both

cat("\nH5 Supported:", H5_supported, "\n")
\end{lstlisting}
\end{rcode}

%==============================================================================
\section{Part 8: Complete Analysis Workflow}
%==============================================================================

\subsection{The 10-Step Process}

Here's your complete workflow from data preparation to hypothesis testing:

\begin{rcode}
\begin{lstlisting}[language=R]
# ========================================
# COMPLETE TAD ANALYSIS WORKFLOW
# ========================================

# STEP 1: Load packages
library(lme4)       # Mixed effects logistic regression
library(survival)   # Cox models
library(ggplot2)    # Plotting
library(dplyr)      # Data manipulation
library(pROC)       # ROC curves
library(metafor)    # Meta-analysis for H5

# STEP 2: Load and prepare data
dat <- read.csv("TAD_data.csv")

# Check structure
str(dat)
summary(dat)

# STEP 3: Standardize all predictors (CRITICAL!)
predictors <- c("livestock_density", "vaccination_coverage",
                "wildlife_proximity", "tree_cover", "water_bodies",
                "wetlands", "temperature", "precipitation",
                "road_density", "culling_coverage", "GDP",
                "distance_to_wildlife")

# Z-score standardization (mean=0, SD=1)
dat[predictors] <- scale(dat[predictors])

# Verify standardization
colMeans(dat[predictors])  # Should be ~0
apply(dat[predictors], 2, sd)  # Should be 1

# STEP 4: Build persistence model
model_persist <- glmer(
  persistence ~ 
    # Main effects
    livestock_density + 
    vaccination_coverage + 
    wildlife_proximity + 
    tree_cover + 
    water_bodies + 
    wetlands +
    temperature + 
    precipitation + 
    road_density + 
    culling_coverage + 
    GDP + 
    wildlife_origin + 
    distance_to_wildlife + 
    # Interactions
    livestock_density:vaccination_coverage +
    wildlife_proximity:culling_coverage +
    wildlife_proximity:tree_cover +
    wetlands:livestock_density +
    road_density:wildlife_proximity +
    # Random effects
    (1|disease/country/admin2),
  family = binomial,
  data = dat,
  control = glmerControl(optimizer = "bobyqa")
)

# Check convergence
if(any(grepl("failed to converge", model_persist@optinfo$conv$lme4$messages))) {
  warning("Model did not converge! Try simplifying or rescaling.")
}

# STEP 5: Model diagnostics
summary(model_persist)

# Check for collinearity
library(car)
vif_values <- vif(model_persist)
print(vif_values)
if(any(vif_values > 5)) {
  warning("High collinearity detected (VIF > 5)")
}

# STEP 6: Build reemergence model
model_reemerge <- coxph(
  Surv(time, event) ~ 
    livestock_density + 
    vaccination_coverage + 
    wildlife_proximity + 
    tree_cover + 
    water_bodies + 
    wetlands +
    temperature + 
    precipitation + 
    road_density + 
    culling_coverage + 
    GDP + 
    wildlife_origin + 
    distance_to_wildlife + 
    livestock_density:vaccination_coverage +
    wildlife_proximity:culling_coverage +
    wildlife_proximity:tree_cover +
    wetlands:livestock_density +
    road_density:wildlife_proximity +
    frailty(disease/country/admin2),
  data = dat
)

# Check proportional hazards assumption
test_ph <- cox.zph(model_reemerge)
print(test_ph)
if(any(test_ph$table[, "p"] < 0.05)) {
  warning("Proportional hazards assumption violated for some predictors")
}

# STEP 7: Extract and organize results
# Persistence
persist_coef <- summary(model_persist)$coefficients
persist_OR <- exp(persist_coef[, "Estimate"])
persist_CI <- exp(confint(model_persist))

results_persist <- data.frame(
  Predictor = rownames(persist_coef),
  Beta = persist_coef[, "Estimate"],
  SE = persist_coef[, "Std. Error"],
  OR = persist_OR,
  CI_lower = persist_CI[, 1],
  CI_upper = persist_CI[, 2],
  p_value = persist_coef[, "Pr(>|z|)"]
)

# Reemergence
reemerge_coef <- summary(model_reemerge)$coefficients
results_reemerge <- data.frame(
  Predictor = rownames(reemerge_coef),
  Beta = reemerge_coef[, "coef"],
  SE = reemerge_coef[, "se(coef)"],
  HR = reemerge_coef[, "exp(coef)"],
  CI_lower = exp(reemerge_coef[, "coef"] - 1.96*reemerge_coef[, "se(coef)"]),
  CI_upper = exp(reemerge_coef[, "coef"] + 1.96*reemerge_coef[, "se(coef)"]),
  p_value = reemerge_coef[, "Pr(>|z|)"]
)

# Round for readability
results_persist[, 2:6] <- round(results_persist[, 2:6], 3)
results_reemerge[, 2:6] <- round(results_reemerge[, 2:6], 3)

# Save results
write.csv(results_persist, "results_persistence.csv", row.names = FALSE)
write.csv(results_reemerge, "results_reemergence.csv", row.names = FALSE)

# STEP 8: Test hypotheses systematically
source("test_hypotheses.R")  # Script with H1-H5 tests from earlier

# STEP 9: Create visualizations
# Forest plot
library(forestplot)

# Prepare data for forest plot
plot_data <- results_persist[!grepl("Intercept|:", results_persist$Predictor), ]

forestplot(
  labeltext = plot_data$Predictor,
  mean = plot_data$OR,
  lower = plot_data$CI_lower,
  upper = plot_data$CI_upper,
  xlab = "Odds Ratio (95% CI)",
  zero = 1,
  col = fpColors(box = "royalblue", line = "darkblue")
)

# Interaction plot (if significant)
# [Code from earlier sections]

# STEP 10: Model performance
# AUC for persistence
pred_prob <- predict(model_persist, type = "response", re.form = NA)
roc_curve <- roc(dat$persistence, pred_prob)
auc_value <- auc(roc_curve)
cat("Persistence model AUC:", round(auc_value, 3), "\n")

# C-index for reemergence
concordance_value <- summary(model_reemerge)$concordance[1]
cat("Reemergence model C-index:", round(concordance_value, 3), "\n")

# DONE!
cat("\n===========================================\n")
cat("Analysis complete! Check results files.\n")
cat("===========================================\n")
\end{lstlisting}
\end{rcode}

%==============================================================================
\section{Part 9: Common Mistakes to Avoid}
%==============================================================================

\begin{warning}
\textbf{Mistake 1: Forgetting to Standardize Predictors}

\textbf{Wrong:}
\begin{lstlisting}[language=R]
model <- glmer(persistence ~ livestock_density + vaccination + ...)
\end{lstlisting}

\textbf{Why it's wrong:} Livestock (range: 100--10,000) and vaccination (range: 0--1) are on different scales. Coefficients are not comparable.

\textbf{Fix:}
\begin{lstlisting}[language=R]
dat[predictors] <- scale(dat[predictors])
model <- glmer(persistence ~ livestock_density + vaccination + ...)
\end{lstlisting}
\end{warning}

\begin{warning}
\textbf{Mistake 2: Interpreting Non-Significant Results}

\textbf{Wrong:} ``Vaccination showed a trend toward reducing persistence (OR = 0.78, $p = 0.12$)''

\textbf{Why it's wrong:} $p = 0.12$ means insufficient evidence. Could be no effect.

\textbf{Fix:} ``Vaccination was not significantly associated with persistence (OR = 0.78, 95\% CI: 0.57--1.07, $p = 0.12$)''
\end{warning}

\begin{warning}
\textbf{Mistake 3: Ignoring Convergence Warnings}

If you see ``failed to converge'', your results are unreliable.

\textbf{Fixes:}
\begin{itemize}
    \item Rescale predictors: \texttt{dat\$livestock <- dat\$livestock / 1000}
    \item Change optimizer: \texttt{control = glmerControl(optimizer = "bobyqa")}
    \item Simplify random effects: Use \texttt{(1|disease) + (1|country)} instead of nested
\end{itemize}
\end{warning}

%==============================================================================
\section{Part 10: Writing Your Results}
%==============================================================================

\subsection{Results Section Template}

\textbf{Introduction paragraph:}
\begin{quote}
We fitted mixed-effects logistic regression models to identify predictors of TAD persistence (n = 206 post-outbreak periods across 15 diseases and 8 countries) and Cox proportional hazards models to analyze time to reemergence (n = 143 disease-free periods). All continuous predictors were z-score standardized prior to analysis. Random intercepts for disease, country, and administrative unit accounted for hierarchical data structure.
\end{quote}

\textbf{Model performance:}
\begin{quote}
The persistence model demonstrated good discrimination (AUC = 0.78, 95\% CI: 0.74--0.82) and adequate calibration. The reemergence model achieved a C-index of 0.74 (95\% CI: 0.69--0.79), indicating moderate predictive ability.
\end{quote}

\textbf{Random effects:}
\begin{quote}
Random effects variance components indicated substantial heterogeneity in baseline risk across diseases ($\sigma^2 = 0.68$), countries ($\sigma^2 = 0.41$), and administrative units ($\sigma^2 = 0.23$), justifying the multilevel modeling approach.
\end{quote}

\textbf{Hypothesis 1 (example):}
\begin{quote}
Consistent with Hypothesis 1, all ecological risk factors showed significant positive associations with persistence (Table 1). Counties with high livestock density (1 SD above mean) had 69\% higher odds of persistence compared to average-density counties (OR = 1.69, 95\% CI: 1.18--2.41, $p = 0.004$). Similarly, wildlife proximity (OR = 1.32, $p = 0.035$), tree cover (OR = 1.48, $p = 0.008$), and wetlands (OR = 1.25, $p = 0.041$) all independently increased persistence risk. These factors also accelerated reemergence, with hazard ratios ranging from 1.28 to 1.54 (all $p < 0.05$; Table 2), fully supporting H1.
\end{quote}

%==============================================================================
\section{Quick Reference: The Cheat Sheet}
%==============================================================================

\subsection{The Four-Step Interpretation (Use Every Time!)}

\begin{enumerate}
    \item \textbf{SIGN:} Positive = increases, Negative = decreases
    \item \textbf{SIGNIFICANT:} $p < 0.05$? If NO, stop here
    \item \textbf{MAGNITUDE:} OR/HR = $e^\beta$, \% = (OR-1)$\times$100\%
    \item \textbf{PRECISION:} Does CI cross 1.0? Narrow = precise
\end{enumerate}

\subsection{OR and HR Rules}

\begin{table}[h]
\centering
\begin{tabular}{@{}cll@{}}
\toprule
\textbf{Value} & \textbf{Meaning} & \textbf{Example} \\
\midrule
= 1.0 & No effect & Predictor doesn't matter \\
$>$ 1.0 & Risk factor & OR = 1.5 (50\% increase) \\
$<$ 1.0 & Protective & OR = 0.7 (30\% decrease) \\
\bottomrule
\end{tabular}
\end{table}

\subsection{Key R Functions}

\begin{table}[h]
\centering
\small
\begin{tabular}{@{}lp{8cm}@{}}
\toprule
\textbf{Function} & \textbf{What It Does} \\
\midrule
\texttt{glmer()} & Fits mixed-effects logistic regression \\
\texttt{coxph()} & Fits Cox proportional hazards model \\
\texttt{fixef()} & Extracts fixed effects ($\beta$ values) \\
\texttt{exp(fixef())} & Converts $\beta$ to OR \\
\texttt{confint()} & Calculates 95\% confidence intervals \\
\texttt{summary()} & Shows full model output \\
\texttt{scale()} & Standardizes predictors (z-scores) \\
\texttt{vif()} & Checks for collinearity \\
\texttt{cox.zph()} & Tests proportional hazards assumption \\
\bottomrule
\end{tabular}
\end{table}

\subsection{Hypothesis Decision Summary}

\begin{table}[h]
\centering
\tiny
\begin{tabular}{@{}lllll@{}}
\toprule
\textbf{Hypothesis} & \textbf{Coefficients} & \textbf{Sign} & \textbf{Threshold} & \textbf{Rule} \\
\midrule
H1 & $\beta_1, \beta_2, \beta_3, \beta_4$ & All $> 0$ & $p < 0.05$ each & ALL significant \\
H2 & $\beta_8$ & $> 0$ & $p < 0.05$ & Significant \\
H3 & $\beta_9, \beta_{10}$ & Both $< 0$ & $p < 0.05$ each & BOTH significant \\
H4 & $\beta_{12}, \beta_{13}, \beta_{16}$ & $+, -, +$ & $p < 0.05$ each & ALL significant \\
H5 & Disease-specific $\beta$s & Varies & Het. $p < 0.05$ & Meta-regression \\
\bottomrule
\end{tabular}
\end{table}

%==============================================================================
\section{Final Checklist Before Analysis}
%==============================================================================

\begin{itemize}[label=$\square$]
    \item Data loaded and checked for missing values
    \item All predictors standardized using \texttt{scale()}
    \item Random effects structure specified correctly: \texttt{(1|disease/country/admin2)}
    \item Both models fitted without convergence warnings
    \item VIF checked for collinearity (all $< 5$)
    \item Proportional hazards assumption checked for Cox model
    \item Results extracted and organized in tables
    \item Four-step interpretation applied to each predictor
    \item Hypothesis testing code run for H1--H5
    \item Visualizations created (forest plots, interactions)
    \item Results section drafted with clear interpretations
\end{itemize}

%==============================================================================
\section{Getting Help}
%==============================================================================

\textbf{If you get stuck:}
\begin{enumerate}
    \item Check the error message carefully
    \item Review the relevant section in this guide
    \item Try the troubleshooting tips in Section 9
    \item Search for the specific error message online
    
\end{enumerate}


\begin{enumerate}
    \item Run a simple example on subset of data
    \item Prepare specific questions: ``I don't understand why...''
    \item Show your attempted interpretation
    \item Bring this guide with notes
\end{enumerate}

\vspace{1cm}
\begin{center}
\rule{0.8\textwidth}{0.4pt}

\textbf{You can do this!}


\rule{0.8\textwidth}{0.4pt}
\end{center}

\end{document}
